\documentclass[10pt,a4paper]{amsart}
\usepackage{amsmath,amssymb,amsthm,graphicx}
\usepackage{verbatim}
\newenvironment{reference}{\comment}{\endcomment}
%\newenvironment{reference}{}{}
\newenvironment{slogan}{\comment}{\endcomment}
\newenvironment{history}{\comment}{\endcomment}

% For commutative diagrams we use Xy-pic
\usepackage[all]{xy}

% We use 2cell for 2-commutative diagrams.
\xyoption{2cell}
% Theorem environments.
%
\theoremstyle{plain}
\newtheorem{theorem}[subsection]{Theorem}
\newtheorem{proposition}[subsection]{Proposition}
\newtheorem{lemma}[subsection]{Lemma}

\theoremstyle{definition}
\newtheorem{definition}[subsection]{Definition}
\newtheorem{example}[subsection]{Example}
\newtheorem{exercise}[subsection]{Exercise}
\newtheorem{situation}[subsection]{Situation}

\theoremstyle{remark}
\newtheorem{remark}[subsection]{Remark}
\newtheorem{remarks}[subsection]{Remarks}

\numberwithin{equation}{subsection}

% Macros
%
\def\lim{\mathop{\mathrm{lim}}\nolimits}
\def\colim{\mathop{\mathrm{colim}}\nolimits}
\def\Spec{\mathop{\mathrm{Spec}}}
\def\Hom{\mathop{\mathrm{Hom}}\nolimits}
\def\Ext{\mathop{\mathrm{Ext}}\nolimits}
\def\SheafHom{\mathop{\mathcal{H}\!\mathit{om}}\nolimits}
\def\SheafExt{\mathop{\mathcal{E}\!\mathit{xt}}\nolimits}
\def\Sch{\mathit{Sch}}
\def\Mor{\mathop{\mathrm{Mor}}\nolimits}
\def\Ob{\mathop{\mathrm{Ob}}\nolimits}
\def\Sh{\mathop{\mathit{Sh}}\nolimits}
\def\NL{\mathop{N\!L}\nolimits}
\def\CH{\mathop{\mathrm{CH}}\nolimits}
\def\proetale{{pro\text{-}\acute{e}tale}}
\def\etale{{\acute{e}tale}}
\def\QCoh{\mathit{QCoh}}
\def\Ker{\mathop{\mathrm{Ker}}}
\def\Im{\mathop{\mathrm{Im}}}
\def\Coker{\mathop{\mathrm{Coker}}}
\def\Coim{\mathop{\mathrm{Coim}}}

% Boxtimes
%
\DeclareMathSymbol{\boxtimes}{\mathbin}{AMSa}{"02}

%
% Macros for moduli stacks/spaces
%
\def\QCohstack{\mathcal{QC}\!\mathit{oh}}
\def\Cohstack{\mathcal{C}\!\mathit{oh}}
\def\Spacesstack{\mathcal{S}\!\mathit{paces}}
\def\Quotfunctor{\mathrm{Quot}}
\def\Hilbfunctor{\mathrm{Hilb}}
\def\Curvesstack{\mathcal{C}\!\mathit{urves}}
\def\Polarizedstack{\mathcal{P}\!\mathit{olarized}}
\def\Complexesstack{\mathcal{C}\!\mathit{omplexes}}
% \Pic is the operator that assigns to X its picard group, usage \Pic(X)
% \Picardstack_{X/B} denotes the Picard stack of X over B
% \Picardfunctor_{X/B} denotes the Picard functor of X over B
\def\Pic{\mathop{\mathrm{Pic}}\nolimits}
\def\Picardstack{\mathcal{P}\!\mathit{ic}}
\def\Picardfunctor{\mathrm{Pic}}
\def\Deformationcategory{\mathcal{D}\!\mathit{ef}}
\usepackage{cleveref}
\crefname{originalreference}{original reference}{original references}
\Crefname{originalreference}{Original reference}{Original references}
\makeatletter
\newlabel{definition-cone}{{014E}{source}{}{sourcepointer-1}{}}
\newlabel{definition-cone@cref}{{[originalreference][0][]014E}{[1][0][]source}{}{}{}}
\makeatother
\crefname{originallemma}{original lemma}{original lemmas}
\Crefname{originallemma}{Original lemma}{Original lemmas}
\makeatletter
\newlabel{homology-lemma-ses-termwise-split-cochain}{{011J}{source}{}{sourcepointer-2}{}}
\newlabel{homology-lemma-ses-termwise-split-cochain@cref}{{[originallemma][0][]011J}{[1][0][]source}{}{}{}}
\makeatother
\crefname{originalreference}{original reference}{original references}
\Crefname{originalreference}{Original reference}{Original references}
\makeatletter
\newlabel{lemma-functorial-cone}{{014F}{source}{}{sourcepointer-3}{}}
\newlabel{lemma-functorial-cone@cref}{{[originalreference][0][]014F}{[1][0][]source}{}{}{}}
\makeatother
\crefname{originalreference}{original reference}{original references}
\Crefname{originalreference}{Original reference}{Original references}
\makeatletter
\newlabel{lemma-make-injective}{{013N}{source}{}{sourcepointer-4}{}}
\newlabel{lemma-make-injective@cref}{{[originalreference][0][]013N}{[1][0][]source}{}{}{}}
\makeatother
\crefname{originalreference}{original reference}{original references}
\Crefname{originalreference}{Original reference}{Original references}
\makeatletter
\newlabel{lemma-third-isomorphism}{{014K}{source}{}{sourcepointer-5}{}}
\newlabel{lemma-third-isomorphism@cref}{{[originalreference][0][]014K}{[1][0][]source}{}{}{}}
\makeatother
\title{Cone triangles of admissible short exact sequences are homotopy equivalent to the quotient, and arbitrary morphisms reduce to admissible monomorphisms}
\author{The Stacks Project Authors}
\date{stacks-project:dga.tex:local-snapshot-20261012; GNU FDL 1.2 or later}
\begin{document}
\maketitle
\noindent\textit{Editorial note.} All delivered mathematical text is the author's own, copied from the pinned Stacks Project chapter file sources/stacks-dga/pinned-original-source.tex; no statement or proof step was rewritten, reordered or removed.\par\medskip
\noindent\textit{Editorial note.} The counted proof is the author's complete argument: the comparison map C(alpha) to M with its homotopy inverse (s, -delta), the differential compatibility via the characterizing property of delta, the displayed homotopy identity with the explicit homotopy h, the cone-map compatibility check p after (s, -delta) = -delta, and the admissible-replacement reduction of part (2) through lemma-make-injective, lemma-functorial-cone and lemma-third-isomorphism are all retained verbatim.\par\medskip
\noindent\textit{Editorial note.} The label lemma-the-same-up-to-isomorphisms names the counted statement, and the proof environment immediately following it in the source is delivered verbatim as its proof.\par\medskip
\noindent\textit{Editorial note.} The counted statement is carried by the author's own lemma environment in the chapter source.\par\medskip
\noindent\textit{Allowed standard prerequisites (editorial).}\begin{itemize}
\item Termwise split sequences and the characterizing differential. Homology, Lemma homology-lemma-ses-termwise-split-cochain. Supplies the characterization of delta as the unique map with d after s minus s after d equal to alpha after delta, which makes the candidate homotopy inverse (s, -delta) compatible with differentials.
\item Admissible replacement of a morphism by a monomorphism. Same chapter, Lemmas lemma-make-injective, lemma-functorial-cone and lemma-third-isomorphism. Reduces part (2) to part (1): after replacing L by the admissible supertilde L, the two cone triangles are isomorphic, so f may be assumed an admissible monomorphism.
\end{itemize}\medskip
\ifdefined\hypertarget\hypertarget{sourcepointer-1}{}\fi
\noindent\textit{Source reference.} Original reference 014E: sources/stacks-dga/pinned-original-source.tex, source lines 619--619. Prerequisite pointer invoked at proof line(s) 1189; it names a standard fact of the same chapter that the counted argument uses, and every step of the counted proof built on it is retained verbatim in the delivered body.\par\medskip
\ifdefined\hypertarget\hypertarget{sourcepointer-2}{}\fi
\noindent\textit{Source reference.} Original lemma 011J: sources/stacks-chow/tags.txt, source lines 1365--1365. Prerequisite pointer invoked at proof line(s) 1115; it names a standard fact of a sibling chapter that the counted argument uses, and every step of the counted proof built on it is retained verbatim in the delivered body.\par\medskip
\ifdefined\hypertarget\hypertarget{sourcepointer-3}{}\fi
\noindent\textit{Source reference.} Original reference 014F: sources/stacks-dga/pinned-original-source.tex, source lines 644--644. Prerequisite pointer invoked at proof line(s) 1202; it names a standard fact of the same chapter that the counted argument uses, and every step of the counted proof built on it is retained verbatim in the delivered body.\par\medskip
\ifdefined\hypertarget\hypertarget{sourcepointer-4}{}\fi
\noindent\textit{Source reference.} Original reference 013N: sources/stacks-dga/pinned-original-source.tex, source lines 777--777. Prerequisite pointer invoked at proof line(s) 1201; it names a standard fact of the same chapter that the counted argument uses, and every step of the counted proof built on it is retained verbatim in the delivered body.\par\medskip
\ifdefined\hypertarget\hypertarget{sourcepointer-5}{}\fi
\noindent\textit{Source reference.} Original reference 014K: sources/stacks-dga/pinned-original-source.tex, source lines 1016--1016. Prerequisite pointer invoked at proof line(s) 1202; it names a standard fact of the same chapter that the counted argument uses, and every step of the counted proof built on it is retained verbatim in the delivered body.\par\medskip
\begin{lemma}
\label{lemma-the-same-up-to-isomorphisms}
Let $(A, \text{d})$ be a differential graded algebra.
\begin{enumerate}
\item Given an admissible short exact sequence
$0 \to K \xrightarrow{\alpha} L \to M \to 0$
of differential graded $A$-modules there exists a homotopy equivalence
$C(\alpha) \to M$ such that the diagram
$$
\xymatrix{
K \ar[r] \ar[d] & L \ar[d] \ar[r] &
C(\alpha) \ar[r]_{-p} \ar[d] & K[1] \ar[d] \\
K \ar[r]^\alpha & L \ar[r]^\beta &
M \ar[r]^\delta & K[1]
}
$$
defines an isomorphism of triangles in $K(\text{Mod}_{(A, \text{d})})$.
\item Given a morphism of complexes $f : K \to L$
there exists an isomorphism of triangles
$$
\xymatrix{
K \ar[r] \ar[d] & \tilde L \ar[d] \ar[r] &
M \ar[r]_{\delta} \ar[d] & K[1] \ar[d] \\
K \ar[r] & L \ar[r] &
C(f) \ar[r]^{-p} & K[1]
}
$$
where the upper triangle is the triangle associated to a
admissible short exact sequence $K \to \tilde L \to M$.
\end{enumerate}
\end{lemma}

\begin{proof}
Proof of (1). We have $C(\alpha) = L \oplus K$ and we simply define
$C(\alpha) \to M$ via the projection onto $L$ followed by $\beta$.
This defines a morphism of differential graded modules because the
compositions $K^{n + 1} \to L^{n + 1} \to M^{n + 1}$ are zero.
Choose splittings $s : M \to L$ and $\pi : L \to K$ with
$\Ker(\pi) = \Im(s)$ and set
$\delta = \pi \circ \text{d}_L \circ s$ as usual.
To get a homotopy inverse we take
$M \to C(\alpha)$ given by $(s , -\delta)$. This is compatible with
differentials because $\delta^n$ can be characterized as the
unique map $M^n \to K^{n + 1}$ such that
$\text{d} \circ s^n - s^{n + 1} \circ \text{d} = \alpha \circ \delta^n$,
see proof of
Homology, Lemma \ref{homology-lemma-ses-termwise-split-cochain}.
The composition $M \to C(f) \to M$ is the identity.
The composition $C(f) \to M \to C(f)$ is equal to the morphism
$$
\left(
\begin{matrix}
s \circ \beta & 0 \\
-\delta \circ \beta & 0
\end{matrix}
\right)
$$
To see that this is homotopic to the identity map
use the homotopy $h : C(\alpha) \to C(\alpha)$
given by the matrix
$$
\left(
\begin{matrix}
0 & 0 \\
\pi & 0
\end{matrix}
\right) :
C(\alpha) = L \oplus K
\to
L \oplus K = C(\alpha)
$$
It is trivial to verify that
$$
\left(
\begin{matrix}
1 & 0 \\
0 & 1
\end{matrix}
\right)
-
\left(
\begin{matrix}
s \\
-\delta
\end{matrix}
\right)
\left(
\begin{matrix}
\beta & 0
\end{matrix}
\right)
=
\left(
\begin{matrix}
\text{d} & \alpha \\
0 & -\text{d}
\end{matrix}
\right)
\left(
\begin{matrix}
0 & 0 \\
\pi & 0
\end{matrix}
\right)
+
\left(
\begin{matrix}
0 & 0 \\
\pi & 0
\end{matrix}
\right)
\left(
\begin{matrix}
\text{d} & \alpha \\
0 & -\text{d}
\end{matrix}
\right)
$$
To finish the proof of (1) we have to show that the morphisms
$-p : C(\alpha) \to K[1]$ (see
Definition \ref{definition-cone})
and $C(\alpha) \to M \to K[1]$ agree up
to homotopy. This is clear from the above. Namely, we can use the homotopy
inverse $(s, -\delta) : M \to C(\alpha)$
and check instead that the two maps
$M \to K[1]$ agree. And note that
$p \circ (s, -\delta) = -\delta$ as desired.

\medskip\noindent
Proof of (2). We let $\tilde f : K \to \tilde L$,
$s : L \to \tilde L$
and $\pi : L \to L$ be as in
Lemma \ref{lemma-make-injective}. By
Lemmas \ref{lemma-functorial-cone} and \ref{lemma-third-isomorphism}
the triangles $(K, L, C(f), i, p)$ and
$(K, \tilde L, C(\tilde f), \tilde i, \tilde p)$
are isomorphic. Note that we can compose isomorphisms of
triangles. Thus we may replace $L$ by
$\tilde L$ and $f$ by $\tilde f$. In other words
we may assume that $f$ is an admissible monomorphism.
In this case the result follows from part (1).
\end{proof}

\end{document}
